\pagestyle{empty}
\tight
\begin{tblr}{width=\textwidth,colspec={|Q[c,m,1.9cm]|X[l,m]|},hlines,vlines,hspan=minimal,row{1}={bg=headblue},rowsep=1pt,colsep=3pt}
\SetCell{c,m}\hfil\logo\hfil & \textbf{POLITEKNIK NEGERI MALANG}\par\textbf{JURUSAN TEKNOLOGI INFORMASI}\par\textbf{PROGRAM STUDI: D4 TEKNIK INFORMATIKA} \\
\SetCell[c=2]{c} \textbf{RENCANA PEMBELAJARAN SEMESTER (RPS)} & \\
\end{tblr}
\begin{longtblr}{width=\textwidth,colspec={|Q[l,m,1.9cm]|X[0.85,l,m]|X[1.45,l,m]|X[0.75,l,m]|X[1.15,l,m]|},hlines,vlines,hspan=minimal,rowsep=1pt,colsep=2pt,row{1,3}={bg=headgray,font=\bfseries}}
MATA KULIAH & KODE & BOBOT (SKS)/jam & SEMESTER & TGL. PENYUSUNAN \\
Big Data & RTI256106 & 2 SKS/4 jam & 6 & 1 Juli 2025 \\
OTORISASI & Dosen Pengembang RPS & Koordinator RMK & \SetCell[c=2]{l} Ka PRODI &  \\
 & Vit Zuraida, S.Kom., M.Kom. &  & \SetCell[c=2]{l}{\vspace{8mm}\par Dr. Ely Setyo Astuti, S.T., M.T.} &  \\
\SetCell[r=12]{l} Capaian Pembelajaran (CP) & \SetCell[c=4]{l,bg=headgray,font=\bfseries} Capaian Pembelajaran Lulusan Program Studi (CPL-Prodi) &  &  &  \\
 & CPL07 & \SetCell[c=3]{l} Mampu mengembangkan sistem komputasi cerdas dan melakukan analisis data berbasis kecerdasan buatan untuk pengambilan keputusan. &  &  \\
 & CPL05 & \SetCell[c=3]{l} Mampu menghasilkan solusi teknologi komputasi dan infrastruktur digital yang sesuai dengan kebutuhan organisasi dan pengguna. &  &  \\
 & CPL09 & \SetCell[c=3]{l} Mampu menjelaskan cara kerja sistem komputer dan menerapkan pendekatan komputasi untuk menyelesaikan masalah. &  &  \\
 & \SetCell[c=4]{l,bg=headgray,font=\bfseries} Capaian Pembelajaran mata kuliah (CPL-MK) &  &  &  \\
 & CPMK0702 & \SetCell[c=3]{l} Mahasiswa mampu memproses data skala besar menggunakan arsitektur paralel dan terdistribusi. &  &  \\
 & CPMK0904 & \SetCell[c=3]{l} Mahasiswa mampu menjelaskan prinsip kerja sistem paralel-terdistribusi pada pemrosesan data skala besar. &  &  \\
 & CPMK0503 & \SetCell[c=3]{l} Mampu merancang dan menerapkan solusi komputasi paralel dan terdistribusi untuk kebutuhan IoT, Big Data, dan cloud. &  &  \\
 & \SetCell[c=4]{l,bg=headgray,font=\bfseries} Kemampuan akhir tiap tahapan belajar (Sub-CPMK) &  &  &  \\
 & SCPMK0702-05001 & \SetCell[c=3]{l} Mahasiswa mampu merancang dan mengimplementasikan pipeline pemrosesan Big Data menggunakan ekosistem Hadoop/Spark. &  &  \\
 & SCPMK0904-05002 & \SetCell[c=3]{l} Mahasiswa mampu menjelaskan arsitektur dan prinsip kerja sistem terdistribusi untuk batch processing dan stream processing. &  &  \\
 & SCPMK0503-05003 & \SetCell[c=3]{l} Mahasiswa mampu merancang dan menerapkan solusi komputasi paralel dan terdistribusi untuk kebutuhan Big Data dan cloud computing. &  &  \\
\SetCell[r=2]{l} Penilaian & \SetCell[c=4]{l} *Nilai CPMK didapat dari agregasi nilai SUB-CPMK yang dibebankan &  &  &  \\
 & \SetCell[c=4]{l}{\tiny\begin{tblr}{width=\linewidth,colspec={|Q[c,m,0.1\linewidth]|Q[c,m,0.16\linewidth]|X[l,m]|X[l,m]|X[l,m]|X[l,m]|X[l,m]|X[l,m]|Q[c,m,0.08\linewidth]|},hlines,vlines,hspan=minimal,rowsep=1pt,colsep=0.7pt,row{1,2}={bg=headgray,font=\bfseries}}
Kode CPMK & Kode SCPMK & \SetCell[c=6]{c} Bobot per Bentuk Penilaian & & & & & & Total Bobot \\
 & & Tugas & Kuis 1 & UTS & Kuis 2 & PBL & UAS & \\
CPMK0702 & SCPMK0702-05001 & 18\%: pipeline Big Data & 0\% & 7\%: implementasi pipeline & 5\%: Spark dan streaming & 17\%: implementasi pipeline & 15\%: demo pipeline & 62\% \\
CPMK0904 & SCPMK0904-05002 & 12\%: arsitektur Big Data & 5\%: konsep Big Data & 8\%: arsitektur distribusi & 0\% & 0\% & 8\%: analisis arsitektur & 33\% \\
CPMK0503 & SCPMK0503-05003 & 0\% & 0\% & 0\% & 0\% & 5\%: solusi komputasi paralel-terdistribusi cloud & 0\% & 5\% \\
\SetCell[c=8]{r}\textbf{Total Per Penilaian} & & & & & & & & \textbf{100\%} \\
\end{tblr}} &  &  &  \\
Diskripsi Singkat Mata Kuliah & \SetCell[c=4]{l} Big Data membangun kemampuan merancang dan mengimplementasikan pipeline pemrosesan data skala besar menggunakan ekosistem Hadoop dan Spark, mencakup batch processing, stream processing, dan visualisasi data. &  &  &  \\
Bahan Kajian / Materi Pembelajaran & \SetCell[c=4]{l}\begin{enumerate}\item Konsep Big Data: 4V (Volume, Velocity, Variety, Veracity), arsitektur Lambda, Kappa, dan HDFS\item Pemrosesan terdistribusi: MapReduce, Hadoop ekosistem (YARN, Hive, Pig), dan Apache Spark (RDD, DataFrame, SQL)\item Stream processing: Spark Streaming, Apache Kafka, dan pemrosesan event real-time\item Data lakehouse: Delta Lake, Apache Iceberg, NoSQL (Cassandra, MongoDB), dan visualisasi Big Data\end{enumerate} &  &  &  \\
\SetCell[r=2]{l} Pustaka & \SetCell[c=4]{l} \textbf{Utama:}\begin{enumerate}\item White, T. 2015. Hadoop: The Definitive Guide. O'Reilly.\item Chambers, B. 2018. Spark: The Definitive Guide. O'Reilly.\item Marz, N. 2015. Big Data: Principles and Best Practices. Manning.\end{enumerate} &  &  &  \\
 & \SetCell[c=4]{l} \textbf{Pendukung:} Materi LMS, dokumentasi resmi Apache Hadoop dan Spark, dan jobsheet praktikum. &  &  &  \\
Media Pembelajaran & \SetCell[c=2]{l} \textbf{Software:} Apache Hadoop, Apache Spark, Apache Kafka, Docker, Jupyter Notebook, LMS. &  & \SetCell[c=2]{l} \textbf{Hardware:} Komputer/laptop, cluster (fisik/cloud), proyektor. &  \\
Nama Dosen Pengampu & \SetCell[c=4]{l} Tim Pengajar Program Studi D4 TI &  &  &  \\
Matakuliah Syarat & \SetCell[c=4]{l} - &  &  &  \\
\end{longtblr}
\begin{longtblr}{width=\textwidth,colspec={|Q[c,m,0.06\textwidth]|X[1.35,l,m]|X[1.08,l,m]|X[1.16,l,m]|X[0.7,l,m]|X[1.24,l,m]|X[1.06,l,m]|X[0.98,l,m]|Q[c,m,0.05\textwidth]|},hlines,vlines,hspan=minimal,rowhead=2,row{1,2}={bg=headgreen,font=\bfseries},rowsep=1pt,colsep=1.5pt}
Mgg.\par Ke & Kemampuan Akhir Yang Direncanakan (Sub-CP-MK) & Bahan kajian (Materi Pembelajaran) & Bentuk dan Metode Pembelajaran & Estimasi Waktu & Pengalaman Belajar Mahasiswa & Kriteria \& Bentuk Penilaian & Indikator Penilaian & Bobot Penilaian (\%) \\
(1) & (2) & (3) & (4) & (5) & (6) & (7) & (8) & (9) \\
1 & SCPMK0904-05002 & Pengantar Big Data: 4V dan paradigma. & Modalitas: Blended Learning. Bentuk: Luring/Daring. Metode: Case Method, praktik Big Data, studi kasus, dan peer review. & 1 x 2 x 50' tatap muka; 1 x 2 x 50' tugas/praktik mandiri. & Mahasiswa mengerjakan aktivitas terarah untuk pengantar big data: 4v dan paradigma dan menunjukkan evidence hasil belajar. & Bentuk penilaian: Pengantar Big Data: 4V dan Paradigma. Mengacu rubrik RTM. & Ketepatan konsep; kualitas implementasi; validasi dan dokumentasi. & 2 \\
2 & SCPMK0904-05002 & Arsitektur Big Data: Lambda dan Kappa. & Modalitas: Blended Learning. Bentuk: Luring/Daring. Metode: Case Method, praktik Big Data, studi kasus, dan peer review. & 1 x 2 x 50' tatap muka; 1 x 2 x 50' tugas/praktik mandiri. & Mahasiswa mengerjakan aktivitas terarah untuk arsitektur big data: lambda dan kappa dan menunjukkan evidence hasil belajar. & Bentuk penilaian: Arsitektur Big Data: Lambda dan Kappa. Mengacu rubrik RTM. & Ketepatan konsep; kualitas implementasi; validasi dan dokumentasi. & 3 \\
3 & SCPMK0904-05002 & HDFS: Distributed File System. & Modalitas: Blended Learning. Bentuk: Luring/Daring. Metode: Case Method, praktik Big Data, studi kasus, dan peer review. & 1 x 2 x 50' tatap muka; 1 x 2 x 50' tugas/praktik mandiri. & Mahasiswa mengerjakan aktivitas terarah untuk hdfs: distributed file system dan menunjukkan evidence hasil belajar. & Bentuk penilaian: HDFS: Distributed File System. Mengacu rubrik RTM. & Ketepatan konsep; kualitas implementasi; validasi dan dokumentasi. & 3 \\
4 & SCPMK0904-05002 & Kuis 1 konsep dan arsitektur Big Data. & Modalitas: Blended Learning. Bentuk: Luring/Daring. Metode: Case Method, praktik Big Data, studi kasus, dan peer review. & 1 x 2 x 50' tatap muka; 1 x 2 x 50' tugas/praktik mandiri. & Mahasiswa mengerjakan aktivitas terarah untuk kuis 1 konsep dan arsitektur big data dan menunjukkan evidence hasil belajar. & Bentuk penilaian: Kuis 1 Konsep dan Arsitektur Big Data. Mengacu rubrik RTM. & Ketepatan konsep; kualitas implementasi; validasi dan dokumentasi. & 5 \\
5 & SCPMK0702-05001 & MapReduce: paradigma pemrosesan terdistribusi. & Modalitas: Blended Learning. Bentuk: Luring/Daring. Metode: Case Method, praktik Big Data, studi kasus, dan peer review. & 1 x 2 x 50' tatap muka; 1 x 2 x 50' tugas/praktik mandiri. & Mahasiswa mengerjakan aktivitas terarah untuk mapreduce: paradigma pemrosesan terdistribusi dan menunjukkan evidence hasil belajar. & Bentuk penilaian: MapReduce: Paradigma Pemrosesan Terdistribusi. Mengacu rubrik RTM. & Ketepatan konsep; kualitas implementasi; validasi dan dokumentasi. & 3 \\
6 & SCPMK0702-05001 & Hadoop ekosistem: YARN, Hive, Pig. & Modalitas: Blended Learning. Bentuk: Luring/Daring. Metode: Case Method, praktik Big Data, studi kasus, dan peer review. & 1 x 2 x 50' tatap muka; 1 x 2 x 50' tugas/praktik mandiri. & Mahasiswa mengerjakan aktivitas terarah untuk hadoop ekosistem: yarn, hive, pig dan menunjukkan evidence hasil belajar. & Bentuk penilaian: Hadoop Ekosistem: YARN, Hive, Pig. Mengacu rubrik RTM. & Ketepatan konsep; kualitas implementasi; validasi dan dokumentasi. & 3 \\
7 & SCPMK0702-05001 & Apache Spark: RDD, DataFrame, Spark SQL. & Modalitas: Blended Learning. Bentuk: Luring/Daring. Metode: Case Method, praktik Big Data, studi kasus, dan peer review. & 1 x 2 x 50' tatap muka; 1 x 2 x 50' tugas/praktik mandiri. & Mahasiswa mengerjakan aktivitas terarah untuk apache spark: rdd, dataframe, spark sql dan menunjukkan evidence hasil belajar. & Bentuk penilaian: Apache Spark: RDD, DataFrame, Spark SQL. Mengacu rubrik RTM. & Ketepatan konsep; kualitas implementasi; validasi dan dokumentasi. & 2 \\
8 & SCPMK0904-05002, SCPMK0702-05001 & UTS arsitektur dan MapReduce. & Modalitas: Blended Learning. Bentuk: Luring/Daring. Metode: Case Method, praktik Big Data, studi kasus, dan peer review. & 1 x 2 x 50' tatap muka; 1 x 2 x 50' tugas/praktik mandiri. & Mahasiswa mengerjakan aktivitas terarah untuk uts arsitektur dan mapreduce dan menunjukkan evidence hasil belajar. & Bentuk penilaian: UTS Arsitektur dan MapReduce. Mengacu rubrik RTM. & Ketepatan konsep; kualitas implementasi; validasi dan dokumentasi. & 15 \\
9 & SCPMK0702-05001 & Spark Streaming dan pemrosesan real-time. & Modalitas: Blended Learning. Bentuk: Luring/Daring. Metode: Case Method, praktik Big Data, studi kasus, dan peer review. & 1 x 2 x 50' tatap muka; 1 x 2 x 50' tugas/praktik mandiri. & Mahasiswa mengerjakan aktivitas terarah untuk spark streaming dan pemrosesan real-time dan menunjukkan evidence hasil belajar. & Bentuk penilaian: Spark Streaming dan Pemrosesan Real-Time. Mengacu rubrik RTM. & Ketepatan konsep; kualitas implementasi; validasi dan dokumentasi. & 3 \\
10 & SCPMK0702-05001 & Apache Kafka: message broker dan event streaming. & Modalitas: Blended Learning. Bentuk: Luring/Daring. Metode: Case Method, praktik Big Data, studi kasus, dan peer review. & 1 x 2 x 50' tatap muka; 1 x 2 x 50' tugas/praktik mandiri. & Mahasiswa mengerjakan aktivitas terarah untuk apache kafka: message broker dan event streaming dan menunjukkan evidence hasil belajar. & Bentuk penilaian: Apache Kafka: Message Broker dan Event Streaming. Mengacu rubrik RTM. & Ketepatan konsep; kualitas implementasi; validasi dan dokumentasi. & 3 \\
11 & SCPMK0702-05001 & Kuis 2 Spark dan streaming. & Modalitas: Blended Learning. Bentuk: Luring/Daring. Metode: Case Method, praktik Big Data, studi kasus, dan peer review. & 1 x 2 x 50' tatap muka; 1 x 2 x 50' tugas/praktik mandiri. & Mahasiswa mengerjakan aktivitas terarah untuk kuis 2 spark dan streaming dan menunjukkan evidence hasil belajar. & Bentuk penilaian: Kuis 2 Spark dan Streaming. Mengacu rubrik RTM. & Ketepatan konsep; kualitas implementasi; validasi dan dokumentasi. & 5 \\
12 & SCPMK0702-05001 & Data lakehouse: Delta Lake, Apache Iceberg. & Modalitas: Blended Learning. Bentuk: Luring/Daring. Metode: Case Method, praktik Big Data, studi kasus, dan peer review. & 1 x 2 x 50' tatap muka; 1 x 2 x 50' tugas/praktik mandiri. & Mahasiswa mengerjakan aktivitas terarah untuk data lakehouse: delta lake, apache iceberg dan menunjukkan evidence hasil belajar. & Bentuk penilaian: Data Lakehouse: Delta Lake, Apache Iceberg. Mengacu rubrik RTM. & Ketepatan konsep; kualitas implementasi; validasi dan dokumentasi. & 2 \\
13 & SCPMK0702-05001 & NoSQL untuk Big Data: Cassandra, MongoDB. & Modalitas: Blended Learning. Bentuk: Luring/Daring. Metode: Case Method, praktik Big Data, studi kasus, dan peer review. & 1 x 2 x 50' tatap muka; 1 x 2 x 50' tugas/praktik mandiri. & Mahasiswa mengerjakan aktivitas terarah untuk nosql untuk big data: cassandra, mongodb dan menunjukkan evidence hasil belajar. & Bentuk penilaian: NoSQL untuk Big Data: Cassandra, MongoDB. Mengacu rubrik RTM. & Ketepatan konsep; kualitas implementasi; validasi dan dokumentasi. & 2 \\
14 & SCPMK0702-05001 & Visualisasi data skala besar. & Modalitas: Blended Learning. Bentuk: Luring/Daring. Metode: Case Method, praktik Big Data, studi kasus, dan peer review. & 1 x 2 x 50' tatap muka; 1 x 2 x 50' tugas/praktik mandiri. & Mahasiswa mengerjakan aktivitas terarah untuk visualisasi data skala besar dan menunjukkan evidence hasil belajar. & Bentuk penilaian: Visualisasi Data Skala Besar. Mengacu rubrik RTM. & Ketepatan konsep; kualitas implementasi; validasi dan dokumentasi. & 2 \\
15 & SCPMK0702-05001 & Studi kasus Big Data industri. & Modalitas: Blended Learning. Bentuk: Luring/Daring. Metode: Case Method, praktik Big Data, studi kasus, dan peer review. & 1 x 2 x 50' tatap muka; 1 x 2 x 50' tugas/praktik mandiri. & Mahasiswa mengerjakan aktivitas terarah untuk studi kasus big data industri dan menunjukkan evidence hasil belajar. & Bentuk penilaian: Studi Kasus Big Data Industri. Mengacu rubrik RTM. & Ketepatan konsep; kualitas implementasi; validasi dan dokumentasi. & 2 \\
16 & SCPMK0702-05001, SCPMK0503-05003 & PBL implementasi pipeline Big Data: solusi komputasi paralel dan terdistribusi berbasis cloud. & Modalitas: Blended Learning. Bentuk: Luring/Daring. Metode: Case Method, praktik Big Data, studi kasus, dan peer review. & 1 x 2 x 50' tatap muka; 1 x 2 x 50' tugas/praktik mandiri. & Mahasiswa mengerjakan aktivitas terarah untuk pbl implementasi pipeline big data, termasuk merancang solusi komputasi paralel-terdistribusi berbasis cloud, dan menunjukkan evidence hasil belajar. & Bentuk penilaian: PBL Implementasi Pipeline Big Data. Mengacu rubrik RTM. & Ketepatan konsep; kualitas implementasi; validasi dan dokumentasi; kesesuaian solusi paralel-terdistribusi berbasis cloud. & 22 \\
17 & SCPMK0702-05001, SCPMK0904-05002 & UAS demo dan analisis pipeline. & Modalitas: Blended Learning. Bentuk: Luring/Daring. Metode: Case Method, praktik Big Data, studi kasus, dan peer review. & 1 x 2 x 50' tatap muka; 1 x 2 x 50' tugas/praktik mandiri. & Mahasiswa mengerjakan aktivitas terarah untuk uas demo dan analisis pipeline dan menunjukkan evidence hasil belajar. & Bentuk penilaian: UAS Demo dan Analisis Pipeline. Mengacu rubrik RTM. & Ketepatan konsep; kualitas implementasi; validasi dan dokumentasi. & 23 \\
\end{longtblr}
